\chapter{Deployment e Integrazione}

\section{Introduzione al Deployment}
La fase finale del progetto SINTON-IA prevede il deployment dei modelli sviluppati e la loro integrazione con la piattaforma preesistente SINTONIA. L'hosting del modulo di Intelligenza Artificiale è stato realizzato sfruttando gli \textit{Space} della piattaforma Hugging Face, permettendo così al sistema principale di interfacciarsi in modo diretto con i modelli predittivi.

\section{Deployment del Modello AI}
Il deployment su Hugging Face è stato strutturato creando un repository dedicato (\texttt{SINTON-IA/sinton-ia-api}) funzionante come un microservizio.

\subsection{Architettura del Microservizio}
L'applicazione è basata su \textbf{FastAPI}, un web framework moderno per la realizzazione di API ad alte prestazioni in Python. Il file principale, \texttt{main.py}, gestisce l'inizializzazione del server e l'esposizione degli endpoint. L'infrastruttura è costruita attorno a un container Docker (definito nel \texttt{Dockerfile}), che permette la riproducibilità dell'ambiente e l'installazione delle dipendenze necessarie dichiarate nel file \texttt{requirements.txt}.

\subsection{Gestione dei Modelli}
Al momento dell'avvio del container, l'applicazione pre-carica in memoria tutti i modelli pre-addestrati e i relativi trasformatori. Questo riduce significativamente la latenza durante le inferenze:
\begin{itemize}
    \item \texttt{final\_model\_depression.pkl}: contiene il modello LightGBM addestrato per la predizione della depressione.
    \item \texttt{logreg\_model.pkl} e \texttt{tfidf\_vectorizer.pkl}: contengono il modello NLP basato su regressione logistica e il vettorizzatore TF-IDF per l'analisi del rischio suicidario (Red Flag).
    \item \texttt{ga\_tuned\_config.json} e \texttt{genetic\_algorithm.py}: includono i parametri ottimali pre-calcolati (Gold Standard generato da Optuna) e la logica sorgente dell'algoritmo genetico adoperato per l'analisi e la gestione della Churn Prevention.
\end{itemize}

\subsection{Endpoint Esposti}
Il modulo AI offre tre principali interfacce RESTful (Endpoint):
\begin{itemize}
    \item \texttt{POST /api/red-flag}: Riceve il testo libero inserito dal paziente e restituisce la probabilità di rischio suicidario, applicando pipeline di pulizia del testo e inferenza con il modello di regressione logistica.
    \item \texttt{POST /api/predict-depression}: Riceve i log giornalieri strutturati del paziente, esegue l'estrazione delle feature necessarie (es. \textit{valence\_trend\_5d} o \textit{max\_neg\_streak}), e restituisce il livello del rischio calcolato col modello LightGBM.
    \item \texttt{POST /api/genetic-algorithm}: Utilizzato per le analisi di Churn Prevention, recepisce i dati sull'aderenza del paziente alla piattaforma e valuta la miglior strategia interattiva facendo uso dell'algoritmo genetico pre-ottimizzato.
\end{itemize}

\section{Integrazione con il Backend SINTONIA}
Il sistema centrale SINTONIA, sviluppato utilizzando il framework \textbf{NestJS}, funge da client logico per il modulo di Intelligenza Artificiale. Il flusso comunicativo avviene attraverso richieste HTTP strutturate e assicurate.



\subsection{Il Servizio di Integrazione \texttt{AiService}}
Il core dell'integrazione è demandato alla classe \texttt{AiService} posta all'interno del backend NestJS. Questo servizio si occupa di interfacciare fluidamente il core di SINTONIA con le risorse remote:

\begin{itemize}
    \item Inizializza i log verificando la correttezza delle configurazioni ambientali per \texttt{AI\_API\_URL}.
    \item Espone il metodo principale \texttt{predict(endpoint, payload)}, che compone la chiamata REST e gestisce efficacemente le risposte o gli errori causati da anomalie di rete (come richieste rifiutate o codice ritornato nei \textit{4xx/5xx}).
    \item Intercetta chiamate specifiche ad altri modelli, come \texttt{predictGeneticAlgorithm()}, in modo da garantire la giusta elaborazione contestuale delle entità inviate per prevenire il drop-out dell'utente.
\end{itemize}

Per mitigare la latenza causata dal \textit{Cold Start} dei servizi serverless su Hugging Face (lo spegnimento automatico dei container inattivi), \texttt{AiService} invia procedure di asincrone di \textit{pre-warming} con la funzione \texttt{wakeUpModels()} per notificare allo \textit{Space} di avviarsi o rimanere in \textit{stand-by} attivo per immediate elaborazioni.

\section{Continuous Integration}
\label{sec:continuous_integration}

La transizione di un modello di Intelligenza Artificiale da un ambiente di ricerca (Jupyter Notebooks) a un ecosistema di produzione clinico richiede standard rigorosi di affidabilità del codice. Per garantire la manutenibilità del progetto SINTON-IA e prevenire il rilascio di codice difettoso, l'architettura è stata dotata di una pipeline di \textbf{Continuous Integration (CI)} completamente automatizzata tramite \textit{GitHub Actions}.

\subsection{Progettazione della Pipeline}
Il workflow di validazione, definito nel file di configurazione \texttt{ci.yml}, si innesca automaticamente per ogni \textit{Push} o \textit{Pull Request}. L'automazione esegue tre macro-attività sequenziali:

\begin{enumerate}
    \item \textbf{Analisi Statica e Linting (Ruff)}: Ruff è stato selezionato come linter per via della sua estrema velocità di esecuzione. La CI interroga l'intera \textit{codebase} per individuare non conformità agli standard PEP8, variabili \textit{undefined} o pattern di codice deprecati, agendo da filtro di primo livello.
    \item \textbf{Validazione dei Notebooks}: Poiché il core analitico risiede nei file \texttt{.ipynb}, la pipeline automatizza il parsing dei JSON sottostanti per verificare che i file non si siano corrotti durante i salvataggi, garantendo la totale replicabilità degli esperimenti scientifici.
    \item \textbf{Integrità del Codice di Modello}: Tutti i file eseguibili del pacchetto \texttt{src} subiscono un check formale di compilazione. Se anche un singolo errore solleva un'eccezione, GitHub interrompe il processo di integrazione, assicurando che solo il codice formalmente corretto e privo di errori di compilazione possa avanzare verso le fasi successive.
\end{enumerate}

L'introduzione di questa architettura CI ha permesso di mantenere elevati standard qualitativi durante tutto lo sviluppo, garantendo che ogni aggiornamento futuro rispetti i requisiti di correttezza necessari all'operatività del sistema.
